\section{External Validation of the Sequencing Approach and Results}
\label{app:external_validation}

The task sequences underlying our empirical analysis are not observed.
O*NET catalogues the tasks that belong to an occupation but not the order in which they are performed, so we
construct that order by prompting a language model (GPT-5-mini).
Every adjacency-based object in the paper rests on that imputation, which raises two concerns.

The first is that the task orderings may be artifacts of how LLMs reason rather than reconstructions of
how work is done.
An LLM asked to arrange a list of tasks may place semantically similar items next to one another, encoding
topical similarity instead of production logic.
Since AI exposure and execution are themselves correlated with what a task is about, such an ordering could manufacture the
contiguity we document without any underlying sequence and contaminate our other exercises.
Appendix~\ref{app:gptPrompts_robustness} showed our orderings are stable across prompts, but stability is
not accuracy, since an instrument can be consistently wrong.

The second is that not all O*NET tasks may be performed as part of a single workflow sequence.
O*NET is a catalogue of what an occupation involves, not a process map, and if its task lists do not describe
work that proceeds in an order, no ordering procedure, however accurate, could deliver the objects our predictions require.
Appendix~\ref{app:frequency_robustness} pruned the O*NET list of tasks to those that are highly frequently performed, which makes
the resulting task sequences more workflow-like, but the caveat still applies.

This appendix addresses both concerns with external data whose work sequences a language model did not impute.
Two corpora serve.
The first is the \emph{Process Classification Framework} (PCF) of the American Productivity \& Quality Center
(APQC), whose member working groups have recorded since 2004 the activities and tasks that make up a firm's
business processes, in the order practitioners perform them \citep{apqc_pcf}.
The second is a collection of process-mining event logs from the 4TU.ResearchData repository \citep{roadfines_log, sepsis_log, bpi2017_log, bpi2020_log}, which
record case by case the activities organizations actually executed and when, so that the order is observed
rather than asserted.

The bulk of the appendix establishes that the LLM can recover an ordering of activities at all, against corpora
whose order is already known.
Appendix~\ref{app:apqc_validation} shuffles the steps of 345 documented process branches and asks the LLM to
order them in a sequence, under the main analysis's prompt and each of its ten alternatives.
It recovers the practitioners' order to a substantial degree, and does so despite the presence of branches whose activities
are mutually exclusive options for which no true order exists.
Appendix~\ref{app:eventlog_validation} repeats the exercise against the event logs, where the data reveal
which precedences are real and the comparison can be restricted to them.
The two benchmarks agree closely about how much sequence structure the instrument recovers.

Subsection~\ref{app:apqc_fragmentation_top} then turns the PCF corpus on the second concern.
It documents how AI exposure and AI execution labels are carried onto PCF steps with natural-language
embeddings, which is the construction behind columns~(4)\textendash(6) of
Table~\ref{tab:fragmentation_index_regression_exposure}, where the PCF estimates of the fragmentation channel
of Prediction~\#3 are reported alongside the main sample and discussed in
Subsection~\ref{sec:fragmentation_prediction}.
It then re-runs Prediction~\#1 on those documented sequences, where AI-executed steps again form chains longer
than chance allows.
We are unable to test Prediction~\#2, however, as the matching procedure is lossy and the surviving sample too small to support
the required comparison.

Taken together, the exercises conducted in this appendix indicate that our findings are artifacts neither of how LLMs
order a list of tasks in a sequence nor of how O*NET catalogues tasks in occupations.

\subsection{Concern \#1: LLM Recovering Real-World Task Orderings}
\label{app:sequencing_validation}

We benchmark the sequencing procedure twice, against two corpora that differ in where their order comes from.
The first records the order practitioners say they work in; the second records the order in which work was
actually carried out.

\subsubsection{Benchmark \#1: Practitioner-Documented Orderings (PCF)}
\label{app:apqc_validation}

We apply our task-ordering approach, exactly as used in the main analysis, to the APQC corpus, and measure how much of the
order practitioners recorded it recovers.
The benchmark is the \emph{Process Classification Framework} (PCF) of APQC in its Cross-Industry version 8.0.
The PCF decomposes work into five nested levels, from 13 Categories through 74 Process Groups, 362 Processes, 1,353 Activities, and 215 Tasks, for 2,017 process elements in total, each with a persistent identifier, a name, and a short description.
What makes it useful here is that within a parent element the children are listed in the order in which they are performed whenever applicable. Ordering steps in ``Order materials and services,'' for instance, runs from \textit{``process/review requisitions''} through \textit{``approve requisitions,''} \textit{``solicit supplier quotes,''} \textit{``create/distribute purchase orders,''} \textit{``expedite orders and satisfy inquiries,''} \textit{``reconcile purchase orders,''} and finally \textit{``research/resolve order exceptions.''}
APQC also publishes industry-specific versions of the framework, which Subsection~\ref{app:apqc_fragmentation_top} draws on; Table~\ref{tab:apqc_process_examples} shows one documented process group from each of two of them.\footnote{Our results are not particular to the cross-industry framework. While there is some degree of heterogeneity in how the LLM recovers sequences from business processes in different industries, the magnitudes of the difference are not large and we do not report them in detail here for brevity.}
\begin{table}[!t]
    \caption{Example Process Groups From Two Industry Frameworks}
    \label{tab:apqc_process_examples}
    \footnotesize
    \centering
    \adjustbox{max width=\textwidth}{\setlength{\tabcolsep}{6pt}
\footnotesize
\begin{tabular}{cll}
\toprule
 & \textbf{Aerospace and Defense} & \textbf{Automotive} \\
Step & Manage product recalls and regulatory audits & Manage parts \\
\midrule
1 & Initiate recall & Manage inventory after sale \\
2 & Assess the likelihood and consequences of occurrence of any hazards & Manage electronic parts catalog \\
3 & Manage recall related communications & Exchange parts/locate parts \\
4 & Submit regulatory reports & Manage returns \\
5 & Monitor and audit recall effectiveness & Rebuild part \\
6 & Manage recall termination & Manage parts retail operations \\
\bottomrule
\multicolumn{3}{l}{\footnotesize Process elements appear in the order APQC's working groups document them.}\\
\end{tabular}
}
    \begin{minipage}{1\linewidth}
\begin{spacing}{0.2}
{\fontsize{9.5pt}{10pt}\selectfont Notes: Each column shows one PCF Process Group with its constituent process elements, in the order APQC documents them. The groups are Hierarchy ID~6.4 of the Aerospace and Defense framework (PCF ID 20110) and Hierarchy ID~6.11 of the Automotive framework (PCF ID 12685).}
\end{spacing}
\end{minipage}\end{table}

We take as our unit of analysis a \emph{process branch}: a parent element together with its ordered children.
Restricting attention to branches with at least three children leaves 345 branches covering 1,917 process elements, with a median of five steps per branch.
Relative to O*NET, the corpus is small and is weighted toward enterprise process work such as finance, procurement, human resources, information technology, and supply chain.

\paragraph{Practitioner orderings are not immutable ground truth.}
\label{app:apqc_validation_caveat}
Before using APQC's listings as a benchmark, we register a caveat about what they are.
The PCF is a documentation convention produced by committees, not a record of observed execution, and for some branches the listed order is one defensible arrangement among several rather than a statement about precedence.

Two examples make the point.
The branch ``Analyze deal options'' lists four children: \textit{``evaluate acquisition options,''} \textit{``evaluate merger options,''} \textit{``evaluate de-merger options,''} and \textit{``evaluate divestiture options.''}
These are mutually exclusive alternatives that a firm considers not necessarily in the presented order.
Reversing the order, from divestiture to acquisition, would be an equally coherent way to write the same branch down.
Similarly, ``Process accounts payable and expense reimbursements'' lists \textit{``process accounts payable,''} then \textit{``process expense reimbursements,''} then \textit{``manage corporate credit cards.''}
These are three distinct disbursement streams that can be run concurrently, and nothing in the process dictates which is documented first.
For branches of this kind any ordering an LLM produces will disagree with APQC roughly half the time by construction, and disagreement carries no information about whether the LLM understands sequence.

We make no attempt to purge such branches from the sample, since doing so would require the very judgment about true sequence that the exercise is meant to test.
We proceed with the full set of branches, which makes the measured overlap a conservative one.

\paragraph{Design.}
\label{app:apqc_validation_design}
For each of the 345 branches we take APQC's published child order as the target, shuffle the children under a fixed seed so that the LLM never observes the published sequence, and ask it to order them.
The prompt is the one used to sequence O*NET tasks in the main analysis (Prompt~\#1 of Appendix~\ref{app:prompts}), carried over verbatim with a single mechanical substitution that maps the occupation domain onto the process domain, replacing ``occupation'' with ``business process'' and ``task'' with ``step.''
We repeat the exercise for each of the ten alternative prompts of Appendix~\ref{app:gptPrompts_robustness}, applying the identical substitution to all eleven so that no prompt is privileged by hand-tuning, and we use the same model and settings as in the main analysis (GPT-5-mini, temperature $0$).

We measure overlap between the LLM's ordering and APQC's with Kendall's $\tau$, computed within each branch and averaged across branches.
We also report the share of APQC-adjacent step pairs that the model places in the correct direction, whose null value is $0.5$, and the share of branches whose full sequence the model reproduces exactly.
Statistical significance is assessed against a permutation null of $1{,}000$ random orderings per branch, which centers $\tau$ at zero.
Of the $3{,}795$ orderings requested, $3{,}793$ returned a complete and parseable permutation of the branch.%\footnote{The two exceptions occur under one alternative prompt, where the model paraphrased a step title rather than returning it verbatim; we require exact matches before accepting an ordering, so those two branches are dropped for that prompt alone.}

\paragraph{Results.}
\label{app:apqc_validation_results}
Table~\ref{tab:apqc_pcf_ordering_validation} reports the main results and Figure~\ref{fig:apqc_tau_by_size} the underlying distributions.
\begin{table}[!t]
    \caption{Overlap Between GPT-Generated and Practitioner-Documented Process Orderings}
    \label{tab:apqc_pcf_ordering_validation}
    \footnotesize
    \centering
    \setlength{\tabcolsep}{8pt}
\begin{tabular}{lc}
\toprule
\addlinespace
\multicolumn{2}{l}{\textit{Mean Kendall $\tau$ against APQC's published process order}} \\
\addlinespace
\quad Main prompt (Prompt 0) & +0.53 \\
\quad Average of 10 alternative prompts & +0.51 \\
\addlinespace
\hline\\[-1.25em]
\multicolumn{2}{l}{\textit{Other overlap measures}} \\
\addlinespace
\quad Adjacent-pair agreement, main prompt (null 0.5) & 0.70 \\
\quad Adjacent-pair agreement, 10-prompt average & 0.70 \\
\quad Exact-order recovery, main prompt & 0.24 \\
\quad Share of branches with $\tau > 0$, main prompt & 0.86 \\
\quad Share of branches with $\tau > 0$, 10-prompt average & 0.90 \\
\hline\\[-1.25em]
Process branches & 345 \\
Process elements ordered & 1,917 \\
\bottomrule
\end{tabular}

    \begin{minipage}{1\linewidth}
\begin{spacing}{0.2}
{\fontsize{9.5pt}{10pt}\selectfont Notes: This table reports the overlap between the orderings GPT-5-mini produces for APQC process branches and the orderings APQC's practitioner working groups published, pooled across all 345 branches.
``Adjacent-pair agreement'' is the share of APQC-adjacent step pairs the model places in the correct direction, with a null value of $0.5$.
``Exact-order recovery'' is the share of branches whose full published sequence the model reproduces.}
\end{spacing}
\end{minipage}\end{table}
\begin{figure}[!t]
\caption{Overlap between GPT-orderings and Original APQC-PCF Order by Branch Length}
\label{fig:apqc_tau_by_size}
\centering
\begin{subfigure}[b]{0.49\textwidth}
    \subcaption{Main prompt}
    \label{fig:apqc_tau_size_main}
    \includegraphics[width=\linewidth]{plots/apqc_pcf_ordering/tau_distribution_by_size_main_share.png}
\end{subfigure}
\hfill
\begin{subfigure}[b]{0.49\textwidth}
    \subcaption{Average of ten alternative prompts}
    \label{fig:apqc_tau_size_alt}
    \includegraphics[width=\linewidth]{plots/apqc_pcf_ordering/tau_distribution_by_size_alt_share.png}
\end{subfigure}
\par\vspace{\fignotesskip}
\begin{minipage}{1\linewidth}
\begin{spacing}{0.2}
{\fontsize{9.5pt}{10pt}\selectfont Notes: Each panel plots the distribution across the 345 APQC process branches of the within-branch Kendall's $\tau$ between GPT-5-mini's ordering and APQC's published ordering, with each bin decomposed by the number of steps in the branch; darker shading denotes longer branches.
Panel~(a) uses the main prompt and panel~(b) averages, within each branch, the $\tau$ obtained under the ten alternative prompts of Appendix~\ref{app:gptPrompts_robustness}.
The red dashed line marks the sample mean and the black dashed line the value a random ordering would deliver in expectation.
Kendall's $\tau$ is coarse on short branches, which is why the extreme bins at both ends are dominated by branches of three and four steps.}
\end{spacing}
\end{minipage}
\end{figure}
Under the main prompt the average Kendall's $\tau$ between the LLM's ordering and APQC's is $0.53$, with a median of $0.60$; $86\%$ of branches have positive $\tau$.
The LLM places $70\%$ of APQC-adjacent step pairs in the correct direction against a null of $50\%$, and reproduces the practitioner sequence \emph{exactly} in $24\%$ of branches.
The last of these is the most demanding, because exact recovery by chance is one in $120$ for a five-step branch and one in $40{,}320$ for an eight-step branch.
Against the permutation null the average $\tau$ corresponds to $z = 22.6$ ($p = 0.001$).
The LLM is not reproducing the order in which the steps were shown to it. The average $\tau$ between its output and the shuffled input list is $0.11$.

Running the ten alternative prompts leaves this essentially unchanged.
Prompt-specific averages span a narrow range, from $0.49$ to $0.53$, so the result is not an artifact of the particular phrasing used in the main prompt.
Averaging the ten alternatives within each branch, which removes the sampling noise of any single formulation, leaves the mean at $0.51$ while raising the share of branches with positive $\tau$ from $86\%$ to $90\%$ and lowering the cross-branch standard deviation from $0.42$ to $0.37$; twelve of the branches that score negative under the main prompt turn positive once averaged.
Prompt choice therefore affects individual branches but not the aggregate picture.

Overlap is higher where the benchmark is a stronger claim about sequence.
Across PCF categories the average $\tau$ ranges from $0.66$ in ``Deliver Services,'' $0.69$ in ``Manage Customer Service,'' and $0.65$ in ``Manage Supply Chain for Physical Products,'' which are transactional and genuinely sequential, down to $0.25$ in ``Develop and Manage Products and Services'' and $0.27$ in ``Manage External Relationships,'' whose children are largely mutually exclusive options of the kind discussed above.

\paragraph{Reading the left tail.}
\label{app:apqc_validation_tail}
A tenth of branches record a negative $\tau$ under the main prompt, and it is worth being clear about what they are, since the sign of $\tau$ overstates the disagreement on short branches.
Kendall's $\tau$ is coarse when there are few steps to order. On a three-step branch it can only take the values $\{-1, -1/3, +1/3, +1\}$, so a single inverted pair already registers as a large negative number, whereas on a ten-step branch the same inversion costs $0.02$ and leaves $\tau$ near $+1$.
Figure~\ref{fig:apqc_tau_by_size} decomposes the distribution by branch length and shows the consequence directly. The mass at $\tau = -1$ consists entirely of three- and four-step branches, and so does most of the spike at $\tau = +1$, where $56\%$ of three-step branches sit against none of the branches with five steps or more.

The tail is not, however, purely an artifact of short branches.
Negative values are about as common among the longest branches as among the shortest, $14\%$ in both cases, and what differs is their magnitude, in that short branches produce reversals while long branches produce values between $-0.03$ and $-0.20$, that is, a handful of misordered pairs within an otherwise correct sequence.
Inspecting the branches that remain negative after averaging over all ten alternative prompts, the recurring case is the one anticipated above: parents whose children are parallel streams or mutually exclusive alternatives, for which no true ordering exists and the sign of $\tau$ is close to a coin flip.

\subsubsection{Benchmark \#2: Observed Process Execution (4TU.ResearchData)}
\label{app:eventlog_validation}

The benchmark above has a limitation it cannot overcome. A documented ordering is an editorial convention, so where a process has no true sequence, disagreement says nothing about the instrument.
Moving from documented to \emph{observed} order removes it.

A process-mining event log records, case by case, the activities an organization actually executed and when.
Two things follow.
First, the ordering is a fact about production rather than an assertion about it.
Second, and more useful here, the data reveal \emph{which precedences are real}. If activity $x$ precedes $y$ in $97\%$ of cases the precedence is genuine, whereas a $50$/$50$ split means the two run concurrently and no ordering could be correct.
We can therefore condition on the pairs that have a true order before asking whether the model recovers it, which is the step the documented benchmark does not permit.

These logs record processes carried out by several people, with the case passing between them, and that property is what makes them the right benchmark here rather than the wrong one.
Read against the short run alone it would be a mismatch, since Section~\ref{sec:shortrun.production} assigns the whole sequence to a single worker.
But the long-run problem of Section~\ref{sec:longrun.production} is about exactly this object: a workflow partitioned into jobs, with a hand-off incurred wherever the work passes from one worker to the next.
What the model holds fixed across both horizons is the order of the steps, not who performs them, so an event log observes the sequence our imputation is meant to recover, at the level of aggregation the long run works with.

\paragraph{Data.}
\label{app:eventlog_data}
We use eight publicly available event logs from the 4TU.ResearchData repository \citep{roadfines_log, sepsis_log, bpi2017_log, bpi2020_log}, listed in Table~\ref{tab:eventlog_overview}: fine collection by a local police force, sepsis care in a hospital emergency department, consumer lending at a financial institution, and five university travel and reimbursement processes.
Together they cover $215{,}715$ cases and $2{,}048{,}949$ recorded events.
Their activity sets are considerably larger than the APQC process branches, between $8$ and $28$ activities against a median of five, so they also test the instrument on long sequences.

From each log we recover one trace per case, the sequence of activity labels in timestamp order, and keep the activities appearing in at least $2\%$ of cases so that the target is not driven by rare exceptions.
We use each activity's \emph{first} occurrence within a case, which handles the loops present in every real log such as repeated laboratory tests or repeated approval rounds.
This yields the two objects we need.
The \emph{observed ordering} ranks activities by their mean normalized first-occurrence position across cases; this is the analogue of APQC's published order.
The \emph{pairwise precedence matrix} records, for every pair, the share of cases containing both in which one precedes the other.
We call a pair \emph{determinate} when one direction holds in at least $80\%$ of cases.
Across the eight logs, $1{,}124$ of $1{,}178$ activity pairs ($95.4\%$) are determinate, so the benchmark is a statement about sequence for almost all of the comparisons we make.

\begin{table}[!t]
    \caption{Event Logs Used for Validation}
    \label{tab:eventlog_overview}
    \begin{center}
    \resizebox{\textwidth}{!}{%
    \begin{tabular}{llrrrr}
    \toprule
    Log & Process & Cases & Events & Activities & Determinate pairs \\
    \midrule
    Road traffic fines & Fine management, local police & 150,370 & 561,470 & 8 & 27 of 28 \\
    Sepsis & Emergency department sepsis care & 1,050 & 15,214 & 15 & 91 of 98 \\
    BPI 2017 & Loan application handling & 31,509 & 1,202,267 & 23 & 238 of 245 \\
    BPI 2020 & Domestic travel declarations & 10,366 & 56,303 & 10 & 41 of 44 \\
    BPI 2020 & International travel declarations & 6,449 & 72,151 & 22 & 218 of 224 \\
    BPI 2020 & Requests for payment & 6,814 & 36,724 & 10 & 39 of 42 \\
    BPI 2020 & Prepaid travel costs & 2,092 & 18,239 & 17 & 128 of 129 \\
    BPI 2020 & Travel permits and declarations & 7,065 & 86,581 & 28 & 342 of 368 \\
    \midrule
    Total & & 215,715 & 2,048,949 & 133 & 1,124 of 1,178 \\
    \bottomrule
    \end{tabular}}
    \end{center}
    \par\vspace{\fignotesskip}
\begin{minipage}{1\linewidth}
\begin{spacing}{0.2}
{\fontsize{9.5pt}{10pt}\selectfont Notes: All logs are openly available from the 4TU.ResearchData repository \citep{bpi2017_log, bpi2020_log, roadfines_log, sepsis_log}.
``Activities'' counts those appearing in at least $2\%$ of a log's cases, which is the set handed to the model.
A pair of activities is ``determinate'' when one of the two orders holds in at least $80\%$ of the cases containing both.}
\end{spacing}
\end{minipage}
\end{table}

\paragraph{Design.}
\label{app:eventlog_design}
The procedure mirrors that of Appendix~\ref{app:apqc_validation}.
For each log we shuffle the activity labels under a fixed seed, so the model never sees the observed sequence, and ask it to order them, supplying a one-line description of the process in place of an occupation title.
We use all eleven prompts, the main specification of Appendix~\ref{app:prompts} and the ten alternatives of Appendix~\ref{app:gptPrompts_robustness}, carried over with the same mechanical substitution that maps the occupation domain onto the process domain.
No log is told which activity starts or ends a case.

We report two levels of analysis.
At the \emph{log level} we compute Kendall's $\tau$ between the model's ordering and the observed ordering, which is directly comparable to the APQC numbers, benchmarked against a permutation null of $1{,}000$ random orderings per log.
At the \emph{pair level} we ask only whether the model places a pair in the direction the traces show.
The pair level is the more informative of the two here, both because it is where the determinacy conditioning bites and because it is closer to what the paper's predictions actually use, since Predictions~\#1 and~\#2 depend on adjacency rather than on the global ordering.

\paragraph{Results.}
\label{app:eventlog_results}
Table~\ref{tab:eventlog_ordering_validation} reports the results and Figure~\ref{fig:eventlog_activity_accuracy} the underlying distribution.

Against a null of $0.5$, the model places $78.8\%$ of all activity pairs in the direction the traces show.
Restricting to the determinate pairs, those the data confirm have a real order, gives $78.9\%$, and weighting each pair by how determinate it is gives $78.8\%$.
That these three numbers coincide is itself informative. The model does no better on the pairs that a real organization orders strictly than on pairs in general, so the aggregate is not being propped up by easy cases.
At the log level the mean Kendall's $\tau$ is $0.48$ with a median of $0.56$, against a permutation null centered at $0.00$ ($z = 6.8$, $p = 0.001$).
The model is not echoing the list it was shown. The average $\tau$ between its output and the shuffled input is $0.007$.

The ten alternative prompts again leave the picture unchanged.
Determinate-pair accuracy under individual prompts spans $0.767$ to $0.817$ and log-level mean $\tau$ spans $0.433$ to $0.559$, so no conclusion here depends on the phrasing used in the main analysis.

The result that matters for the paper is the comparison with Appendix~\ref{app:apqc_validation}.
Restricting both exercises to \emph{adjacent} pairs, the statistic they share, the documented benchmark gives $0.700$ under the main prompt and $0.697$ averaged over the alternatives, while the observed benchmark gives $0.679$ and $0.699$.
Two independent benchmarks, one written down by practitioners and one recovered from what organizations actually did, agree to within roughly two percentage points about how much sequence structure the instrument recovers.

\begin{table}[!t]
    \caption{Overlap Between GPT-Generated Orderings and Observed Process Execution}
    \label{tab:eventlog_ordering_validation}
    \begin{center}
    \input{tables/eventlog_ordering_validation.tex}
    \end{center}
    \par\vspace{\fignotesskip}
\begin{minipage}{1\linewidth}
\begin{spacing}{0.2}
{\fontsize{9.5pt}{10pt}\selectfont Notes: Pairwise precedence accuracy is the share of activity pairs the model places in the direction the event traces show, with a null value of $0.5$.
``Determinate pairs'' are those for which one order holds in at least $80\%$ of the cases containing both activities.
``Adjacent pairs'' are consecutive in the observed ordering and are the statistic directly comparable to the adjacent-pair agreement of Table~\ref{tab:apqc_pcf_ordering_validation}.
Column~(2) pools the ten alternative prompts of Appendix~\ref{app:gptPrompts_robustness}, so its pair counts are ten times those of column~(1).
The random-ordering row is the mean of $1{,}000$ random permutations per log.}
\end{spacing}
\end{minipage}
\end{table}

\begin{figure}[!t]
\caption{Distribution of Overlap With Observed Execution Across Activities}
\label{fig:eventlog_activity_accuracy}
\begin{center}
\begin{subfigure}[b]{0.49\textwidth}
    \subcaption{Main prompt}
    \label{fig:eventlog_acc_main}
    \includegraphics[width=\linewidth]{plots/eventlog_ordering/activity_accuracy_main_share.png}
\end{subfigure}
\hfill
\begin{subfigure}[b]{0.49\textwidth}
    \subcaption{Average of ten alternative prompts}
    \label{fig:eventlog_acc_alt}
    \includegraphics[width=\linewidth]{plots/eventlog_ordering/activity_accuracy_alt_share.png}
\end{subfigure}
\end{center}
\par\vspace{\fignotesskip}
\begin{minipage}{1\linewidth}
\begin{spacing}{0.2}
{\fontsize{9.5pt}{10pt}\selectfont Notes: The unit is an activity ($133$ across the eight logs).
For each activity we take the determinate pairs it belongs to and plot the share the model ordered as the traces show.
The black dashed line marks the coin flip and the red dashed line the mean.
Averaging over prompts raises the share of activities above the coin flip from $86.5\%$ to $91.7\%$ while leaving the mean almost unchanged, the same pattern as in Figure~\ref{fig:apqc_tau_by_size}.
Because activities within a log share one ordering and each pair contributes to two activities, this distribution is descriptive; the inferential statistics are the pooled pair-level accuracy and the log-level permutation test.}
\end{spacing}
\end{minipage}
\end{figure}

\paragraph{Where the instrument fails, and why it understates its own accuracy.}
\label{app:eventlog_failures}
Performance varies sharply across logs, and the variation is diagnostic.
Determinate-pair accuracy is $0.96$ for road traffic fines, $0.95$ for sepsis care, and $0.91$ for international declarations, but $0.56$ for requests for payment and $0.49$ for domestic declarations.
Inspecting the errors, the failures are overwhelmingly about \emph{reading the labels}, not about understanding sequence.

Three examples carry most of the damage.
In the domestic declaration log the model ranks \texttt{Request Payment} first, reading it as the employee's initiating request; in the data it is the second-to-last step, the finance system requesting disbursement.
In the loan log the model treats \texttt{A\_Accepted} as terminal, whereas in the data it is an early state meaning the application entered processing, preceding validation, offer creation, and completion; seven of the twelve most determinate errors are this single activity.
In the sepsis log the model places \texttt{Return ER} mid-visit, while in the data it denotes a patient returning to the emergency department after discharge and is genuinely last.

This matters for how the number should be read.
Activity labels in an event log are the strings an information system emits, not descriptions of work.
Some are terse (\texttt{A\_Concept}), and many encode the actor rather than the action (\texttt{Declaration APPROVED by ADMINISTRATION}).
The model is therefore working with substantially less context than an O*NET task statement provides, and the errors we observe are of a kind that richer task descriptions would largely prevent.
Accuracy here is, in that specific sense, a conservative estimate of what the procedure achieves on the task statements the paper actually uses.


\subsection{Concern \#2: Transferring the Labels onto Documented Sequences}
\label{app:apqc_fragmentation_top}

Estimating anything on APQC sequences requires transferring AI exposure and execution labels, since those are measured over O*NET tasks and have no counterpart defined on PCF elements.
We pool the Cross-Industry framework with APQC's seventeen industry-specific frameworks, discard duplicate
process groups, embed each leaf element and each labeled O*NET task with the language embedding model \texttt{all-mpnet-base-v2}, and match every step to its nearest task by cosine similarity.
Matching process steps to task statements is lossy, so we set a similarity threshold of $0.71$ and carry a
label across only when the match clears it.\footnote{We chose the $0.71$ similarity threshold after manual inspection of some of the matches. The paragraph on the similarity floor below reports how the results move as it is varied.}
A step whose best match falls below the threshold is not discarded.
Doing so would shorten the very sequences the exercise is about, and the shortening would not be random, since
the steps that fail to match are concentrated in the operational and industry-specific parts of the frameworks.
We instead keep every step in place and code the sub-threshold ones as neither AI-exposed nor AI-executed,
which preserves the documented ordering intact and treats an unverifiable label as an absent one.
Process groups with fewer than five steps are dropped, leaving $13{,}482$ steps in $525$ process groups and
$26$ steps per group, against $20$ tasks per occupation in the main sample.

\paragraph{What survives the threshold.}
The $0.71$ cosine similarity threshold is demanding, and it should be read as such before any estimate is.
Only $2{,}067$ steps, $15\%$ of the corpus, clear it, at a mean cosine similarity of $0.75$.
Because every step below it is coded manual and unexposed, the resulting sample is very thin in AI content, where
$12\%$ of steps are AI-exposed and $4\%$ are AI-executed, against $44\%$ and $13\%$ in the O*NET task
universe.
Anything estimated on these sequences therefore asks whether a prediction holds on a corpus carrying roughly a
third of the main sample's AI density, with every label that does survive verified at high similarity.
The caveat applies to the fragmentation estimates in those columns as much as to the chain-length test below, though Subsection~\ref{sec:fragmentation_prediction} explains why the same sparsity works in the fragmentation test's favor.

\paragraph{How much the similarity floor matters.}
The floor is a judgement call, so we vary it from $0.65$ to $0.75$ in steps of $0.01$ and re-run both exercises at each value.
The $525$ process groups are held fixed throughout, since the five-step minimum applies to the raw step count rather than to the labels, so raising the floor only relabels steps as unexposed and unexecuted.
Figure~\ref{fig:apqc_similarity_threshold} reports the result.
Both signs survive the whole range.
The average AI chain exceeds its reshuffle null at all eleven floors, and the fragmentation coefficient is negative at all eleven under each of the three fixed-effect specifications, reaching the $5\%$ level in eight of them without fixed effects.
The magnitude is not monotone, however, running from $-0.26$ at $0.65$ to $-0.63$ at $0.69$ and back to $-0.11$ at $0.72$, and the standard error roughly doubles across the range as the labels thin out.
Neighboring floors can therefore fall on opposite sides of a significance threshold without the underlying estimates differing, as at $0.72$ and $0.73$, whose point estimates each lie inside the other's confidence interval, so the range matters more than any single cell.
We therefore choose the floor on the density of the transferred labels rather than on the estimates it produces.
Raising it thins both the AI exposure and the AI execution content of the corpus, and exposure, the more common of the two labels, is the binding one.
At $0.71$ the corpus retains $12\%$ of its steps as AI-exposed and $4\%$ as AI-executed, and every floor above it drops the exposure share below $10\%$, which is why we settled there.
\begin{figure}[!t]
\caption{PCF Results as the Label-Transfer Similarity Floor Varies}
\label{fig:apqc_similarity_threshold}
\centering
\includegraphics[width=\textwidth]{plots/apqc_similarity_threshold/apqc_similarity_threshold.png}
\par\vspace{\fignotesskip}
\begin{minipage}{1\linewidth}
\begin{spacing}{0.2}
{\fontsize{9.5pt}{10pt}\selectfont Notes: Each point re-runs the label transfer at the indicated cosine similarity floor and re-estimates on the resulting labels.
The $525$ process groups are the same at every floor, since the five-step minimum is applied to the raw step count and not to the labels.
The left panel reports the average AI chain length against the mean of $1{,}000$ within-group step-order reshuffles, with the shaded band covering two standard deviations of that null and the annotation giving the $z$-statistic.
The right panel reports the coefficient on the empirical fragmentation index from Equation~\eqref{eq:fragmentation_index_regression} with $90\%$ confidence intervals, under each of the three fixed-effect specifications of columns~(4)\textendash(6) of Table~\ref{tab:fragmentation_index_regression_exposure}.
The dotted line marks the $0.71$ floor used throughout the paper.}
\end{spacing}
\end{minipage}
\end{figure}

\paragraph{Prediction \#1.}
The average AI chain in the PCF dataset spans $1.16$ steps, against $1.45$ in the main text for O*NET.
The magnitude is small, and mechanically so, because with only $4\%$ of steps AI-executed, there is little
opportunity for long runs to form, and the ceiling on the statistic is far below the main sample's.
What the exercise can still ask is whether the runs that do form are longer than the same labels arranged at
random along the same sequences, and they clearly are.
Benchmarked as in Section~\ref{sec:chainLength_prediction} against $1{,}000$ draws of a within-group
step-order reshuffle, the null averages $1.09$ with a standard deviation of $0.01$, and the observed value
sits at the 100th percentile ($z = 6.2$); against a within-category reassignment that preserves group size,
the null averages $1.06$ and the observed value again sits at the 100th percentile ($z = 10.2$).
The economic magnitude of AI chain length is therefore modest, but the co-occurrence of AI-executed content behind it is unambiguous.
Since no language model placed these steps in order, the contiguity is a property of the documented workflow
rather than of the ordering procedure.

\paragraph{Prediction \#2.}
We do not attempt a counterpart of the neighbor-execution regression, since that test requires the same step
to appear in several workflows with \emph{different} execution outcomes, while transferred labels are a
deterministic function of the matched task.
Restricting to steps that clear the threshold and whose matched task maps to a single DWA appearing in more
than one process group leaves $981$ observations across $137$ DWAs, of which only $9$ contain both an
AI-executed and a non-executed step, against $10{,}708$ tasks across $1{,}748$ DWAs in the baseline specification of
Table~\ref{tab:DWA_regression_aiExecution_mainSample}, so the matched corpus is underpowered for that exercise
by roughly an order of magnitude.

\paragraph{Reading the PCF fragmentation coefficient.}
The fragmentation regression of Equation~\eqref{eq:fragmentation_index_regression} is estimated on this corpus in columns~(4)\textendash(6) of Table~\ref{tab:fragmentation_index_regression_exposure} and discussed in Subsection~\ref{sec:fragmentation_prediction}, which also explains why the sparsity of the transferred labels leaves the fragmentation coefficient identified here and not on the main sample.
One feature of the transfer belongs here rather than in the main text.
All variables are z-scored against their mean and standard deviation within the PCF sample, and the index varies far less across PCF process groups than across O*NET occupations, with a standard deviation of $0.06$ against $0.26$.
A one-standard-deviation move is thus a much smaller move in the index here, and the slope per unit of the index is correspondingly steeper.

\subsection{Summary and Takeaways}
\label{app:external_validation_discussion}

The benchmarks answer the objection that motivated them.
An instrument that merely clustered semantically similar items would have no reason to place seven in ten
consecutive step pairs in the practitioners' direction, or to reproduce a full documented sequence one time in
four, on a corpus that working groups wrote down with no language model involved at any stage.
Nor would it place four in five activity pairs in the direction real organizations executed them, on pairs the
traces confirm have a true order.
That two independent benchmarks, one documented and one observed, agree to within roughly two percentage
points on the statistic they share is the strongest evidence here that what the instrument recovers is
sequence rather than topic.

The overlap is nonetheless imperfect, and two considerations bound how it should be read.
On one side, some of the residual disagreement is indeterminacy in the benchmark rather than error in the
language model's logic, so the measured overlap understates agreement with genuine workflow order.
On the other, the shortfall is not merely noise.
Across the same branches the model's orderings agree with \emph{each other} at an average pairwise $\tau$ of
$0.79$, well above the $0.53$ they achieve against APQC, so the instrument is considerably more
self-consistent than it is accurate.

Three limitations qualify what these exercises establish.
First, they validate the ordering \emph{procedure} rather than the specific O*NET orderings used in the paper,
since neither the PCF nor the event logs is crosswalked to O*NET tasks.
Second, the PCF is not a random sample of work, since it is weighted toward enterprise process work such as finance,
procurement, human resources, information technology, and supply chain.
Third, the PCF serves as a validation sample rather than as the main sample because our AI exposure labels
come from \citet{eloundou2023gpts} and our AI execution labels from the Anthropic Economic Index
\citep{handa2025economictasksperformedai}, both defined over O*NET tasks.
There is no corresponding measure for PCF process elements, and constructing one would mean generating the
very labels the main analysis takes from external sources.
Working with O*NET keeps our measures comparable to the rest of the literature that uses these now-standard
objects, at the cost of not observing workflow order directly; this appendix is our attempt to bound that
cost.
